% =============================================================
%  期末总复习
% =============================================================
\Chapter{终局\quad 考前三十天}
\noindent{\small\kaishu\color{tipc}（课本总复习 P98--P103）}

\begin{blunder}
豆角爸爸小时候考砸了一次，回家大人说：“没关系。”
正准备松口气，大人从柜子里拿出一套卷子：“趁热打铁，做做。”
——那是豆角爸爸第一次知道什么叫“笑里藏刀”。
所以豆角爸爸这份东西里，\textbf{考前一天只做四件事，一件都不是做新题}。
\end{blunder}

\begin{gamble}{赌一把：模拟卷换钓鱼}
\begin{itemize}
  \item 开考前跟小豆角约好一个分数线（\textbf{别定 100，定 85 就好}）。
  \item 达到了，\textbf{换一次钓鱼机会}；没达到，不罚，但要把错题讲给豆角爸爸听。
  \item \textbf{讲得清就再给一次机会}——因为“讲明白”比“考高分”值钱。
\end{itemize}
\end{gamble}

\begin{bonus}
\begin{itemize}
  \item \textbf{角色反转}：让小豆角当老师，给豆角爸爸出十道题，豆角爸爸答。
        \textbf{豆角爸爸故意错三道}，小豆角必须指出错在哪——小豆角指得出来，说明小豆角真会了。
  \item \textbf{错题拍卖}：把错题本里最难的三道题拿出来，讲对的人赢一颗星。
  \item \textbf{考前“免做题卡”}：连续三天抽查全对，换一张“免做一道题”的卡，
        期末前用掉。小豆角会为了这张卡认真起来。
\end{itemize}
\end{bonus}

\begin{unitinfo}
别等考前一星期才开动。八个单元加两次实践，考前一周搞不完，只能搞得小豆角很慌。\\
\textbf{三轮，每天二十分钟，三周够用}：\\
\textbf{第一轮}——回到八个单元，把每章末尾的“5 分钟快检”重做一遍，找出弱的那个。\\
\textbf{第二轮}——按专项打：计算、口诀、图形测量、应用方向，各来一遍。\\
\textbf{第三轮}——做两套模拟卷，练手感和时间。
\end{unitinfo}

\subsection{三轮怎么排}

\begin{center}
\small
\begin{tabularx}{\linewidth}{@{}p{1.8cm}p{2.6cm}Xp{2.8cm}@{}}
\toprule
\textbf{轮次} & \textbf{什么时候} & \textbf{干什么} & \textbf{每天多久} \\
\midrule
第一轮 & 考前第 3 周前 5 天 & 八章的快检各重做一遍，错的记进错题本，标出最弱的章 &
       20 分钟 \\
第二轮 & 考前第 2 周 & 四个专项各练一次，题在下面 & 20--25 分钟 \\
第三轮 & 考前最后一周 & 模拟卷 A、B 各做一次，讲评一次，最后一天只看错题不做新题 &
       30 分钟 \\
考前那天 & 就这一天 & 只看错题本、背口诀、检查文具，别的都别干 & 15 分钟 \\
\bottomrule
\end{tabularx}
\end{center}

\subsection{第一轮：八个单元过一遍}

做完一章打个钩，错的题号写下来。\textbf{错三道以上的那章，第二轮要重点打。}

\begin{center}
\small
\begin{tabularx}{\linewidth}{@{}p{0.8cm}p{5.4cm}p{1.6cm}X@{}}
\toprule
\textbf{序} & \textbf{哪一章} & \textbf{做完了} & \textbf{错题号 / 备注} \\
\midrule
1 & 100 以内加与减（二） & \cb & \blank[4.4cm] \\
2 & 测量（一） & \cb & \blank[4.4cm] \\
3 & 数一数与乘法 & \cb & \blank[4.4cm] \\
4 & 乘法口诀（一） & \cb & \blank[4.4cm] \\
5 & 分一分与除法 & \cb & \blank[4.4cm] \\
6 & 图形的运动（一） & \cb & \blank[4.4cm] \\
7 & 乘法口诀（二） & \cb & \blank[4.4cm] \\
8 & 乘除法的应用（一） & \cb & \blank[4.4cm] \\
\bottomrule
\end{tabularx}
\end{center}

\subsection{第二轮：四个专项}

\subsubsection{专项一：100 以内加减（含竖式）}
\begin{enumerate}[label=\arabic*.,itemsep=3pt]
  \item $46+38=$ \quad 2. $27+65=$ \quad 3. $9+83=$ \quad 4. $52-27=$ \quad
  5. $80-46=$
  \item $100-58=$ \quad 7. $73-38=$ \quad 8. $35+49=$ \quad 9. $60-24=$ \quad
  10. $91-46=$
\end{enumerate}
\textbf{答案}：1. 84\quad 2. 92\quad 3. 92\quad 4. 25\quad 5. 34\quad
6. 42\quad 7. 35\quad 8. 84\quad 9. 36\quad 10. 45\\
\textbf{要求}：第 6--10 题\textbf{必须写竖式}，进位小 1 和退位点都得有。

\subsubsection{专项二：口诀与乘除计算}
\begin{enumerate}[label=\arabic*.,itemsep=3pt]
  \item $7\times8=$ \quad 2. $6\times9=$ \quad 3. $9\times9=$ \quad 4. $8\times8=$ \quad
  5. $7\times7=$
  \item $54\div6=$ \quad 7. $72\div8=$ \quad 8. $63\div7=$ \quad 9. $45\div5=$ \quad
  10. $36\div9=$
  \item $8\times(\ )<60$，最大填 \blank[1.2cm]。\quad
  12. $(\ )\times7<50$，最大填 \blank[1.2cm]。
\end{enumerate}
\textbf{答案}：1. 56\quad 2. 54\quad 3. 81\quad 4. 64\quad 5. 49\quad
6. 9\quad 7. 9\quad 8. 9\quad 9. 9\quad 10. 4\quad 11. 7\quad 12. 7

\subsubsection{专项三：测量与图形}
\begin{enumerate}[label=\arabic*.,itemsep=3pt]
  \item 填单位：教室长约 8 \blank[1.4cm]；一支新铅笔长约 18 \blank[1.4cm]。
  \item $1\ \text{米}=$ \blank[1.2cm] 厘米；\ $250\ \text{厘米}=$ \blank[1.2cm] 米
        \blank[1.2cm] 厘米。
  \item 一条线段从 3 画到 11，长 \blank[1.2cm] 厘米。
  \item 用坏尺子画 7 厘米，从 5 画到 \blank[1.2cm]。
  \item 下面哪些是轴对称图形？（圆、正方形、平行四边形、等腰三角形、等腰梯形）
  \item 陀螺的轴扎在纸板哪儿转得最稳？
  \item 华容道里的人物，每次只能往哪儿移？
\end{enumerate}
\textbf{答案}：1. 米；厘米\quad 2. 100；2 米 50 厘米\quad 3. 8\quad 4. 12\quad
5. 圆、正方形、等腰三角形、等腰梯形（\textbf{平行四边形不是}）\quad
6. 中心\quad 7. 只能横着或竖着。

\subsubsection{专项四：应用与方向}
\begin{enumerate}[label=\arabic*.,itemsep=3pt]
  \item 每盘 6 个苹果，7 盘一共多少个？
  \item 42 个苹果，每 6 个装一盘，能装几盘？
  \item 小明有 8 张卡片，小红有 32 张，小红是小明的几倍？
  \item 一支笔 5 元，买 7 支多少钱？给 50 元找回多少？
  \item 一本书 9 元，买 6 本，60 元够吗？
  \item 小红面向北，她后面是 \blank[1.2cm]，左面是 \blank[1.2cm]。
  \item 地图上一般是上 \blank[1cm] 下 \blank[1cm] 左 \blank[1cm] 右 \blank[1cm]。
\end{enumerate}
\textbf{答案}：1. $6\times7=42$（个）\quad 2. $42\div6=7$（盘）\quad
3. $32\div8=4$（倍）\quad 4. $5\times7=35$（元），$50-35=15$（元）\quad
5. $9\times6=54$（元），$54<60$ 够\quad 6. 南；西\quad 7. 北、南、西、东。

\subsection{第三轮：模拟卷 A（数与代数）}

\begin{tipbox}
\textbf{做题要求}：一次做完，限时 40 分钟；不会的先跳过，做完再回头；
\textbf{做完小豆角自己先检查一遍}，再给豆角爸爸看。
\end{tipbox}

\noindent\textbf{一、直接写出得数（每题 1 分，共 10 分）}
\begin{enumerate}[label=\arabic*.,itemsep=2pt]
  \item $8\times7=$ \blank[1.4cm] \quad 2. $45\div9=$ \blank[1.4cm] \quad
  3. $36+29=$ \blank[1.4cm] \quad 4. $70-26=$ \blank[1.4cm] \quad
  5. $6\times9=$ \blank[1.4cm]
  \item $54\div6=$ \blank[1.4cm] \quad 7. $100-48=$ \blank[1.4cm] \quad
  8. $7\times8=$ \blank[1.4cm] \quad 9. $24\div4=$ \blank[1.4cm] \quad
  10. $56\div7=$ \blank[1.4cm]
\end{enumerate}

\noindent\textbf{二、用竖式计算（每题 3 分，共 12 分）}
\begin{enumerate}[label=\arabic*.,itemsep=6pt]
  \item $38+27=$ \wls[2]
  \item $100-36=$ \wls[2]
  \item $47+29=$ \wls[2]
  \item $82-45=$ \wls[2]
\end{enumerate}

\noindent\textbf{三、填一填（每空 1 分，共 20 分）}
\begin{enumerate}[label=\arabic*.,itemsep=3pt]
  \item 5 个 6 相加是 \blank[1.4cm]，用乘法算式表示是 \blank[3cm]。
  \item $1\ \text{米}=$ \blank[1.2cm] 厘米；\ $300\ \text{厘米}=$ \blank[1.2cm] 米。
  \item 一支铅笔长约 18 \blank[1.4cm]；教室长约 8 \blank[1.4cm]。
  \item 一个数除以 7 得 8，这个数是 \blank[1.4cm]。
  \item $8\times(\ )<60$，括号里最大填 \blank[1.4cm]。
  \item 12 个苹果平均分给 4 人，每人 \blank[1.2cm] 个；
        如果每人 3 个，可以分给 \blank[1.2cm] 人。
  \item $\triangle\triangle\triangle\triangle$ \ $\triangle\triangle\triangle\triangle$ \
        $\triangle\triangle\triangle\triangle$ 表示 \blank[1cm] 个 \blank[1cm]，
        乘法算式是 \blank[2.6cm]。
  \item 在 $6\times7$ 和 $7\times6$ 中间填 $>$、$<$ 或 $=$：\blank[1.4cm]。
  \item 45 是 9 的 \blank[1.2cm] 倍。
  \item 钟面上有 12 个数，每 3 个数一组，可以分成 \blank[1.2cm] 组。
\end{enumerate}

\noindent\textbf{四、判断（对的打“√”，错的打“×”，每题 1 分，共 5 分）}
\begin{enumerate}[label=\arabic*.,itemsep=2pt]
  \item $3+4$ 可以写成 $3\times4$。\hfill（\quad）
  \item 所有的乘法口诀都能写出两个乘法算式。\hfill（\quad）
  \item 1 米比 100 厘米长。\hfill（\quad）
  \item 求 48 是 6 的几倍，用 $48\div6$。\hfill（\quad）
  \item $100-48=62$。\hfill（\quad）
\end{enumerate}

\noindent\textbf{五、选择（每题 1 分，共 5 分）}
\begin{enumerate}[label=\arabic*.,itemsep=2pt]
  \item 下面哪个数是 7 的倍数？\quad A. 42 \quad B. 45 \quad C. 50 \hfill（\quad）
  \item 小红面向北，她右手边是\quad A. 东 \quad B. 南 \quad C. 西 \hfill（\quad）
  \item 一支笔 3 元，买 6 支要\quad A. 9 元 \quad B. 18 元 \quad C. 2 元 \hfill（\quad）
  \item $25\div5=$\quad A. 4 \quad B. 5 \quad C. 6 \hfill（\quad）
  \item 一根绳子长 2 米，剪去 60 厘米，还剩\quad
        A. 140 厘米 \quad B. 130 厘米 \quad C. 60 厘米 \hfill（\quad）
\end{enumerate}

\noindent\textbf{六、解决问题（每题 7 分，共 28 分）}
\begin{enumerate}[label=\arabic*.,itemsep=8pt]
  \item 每盒装 8 个鸡蛋，买 5 盒一共多少个？
        \wls[2]
  \item 有 36 个苹果，每 4 个装一袋，能装几袋？
        \wls[2]
  \item 小明今年 7 岁，爷爷 63 岁，爷爷的年龄是小明的几倍？
        \wls[2]
  \item 一本书 9 元，小红有 50 元，买 5 本够吗？还剩多少元？
        \wls[3]
\end{enumerate}

\begin{ansbox}
\textbf{卷 A 答案}\\
\textbf{一}：1. 56\quad 2. 5\quad 3. 65\quad 4. 44\quad 5. 54\quad
6. 9\quad 7. 52\quad 8. 56\quad 9. 6\quad 10. 8 \\
\textbf{二}：1. 65\quad 2. 64\quad 3. 76\quad 4. 37 \\
\textbf{三}：1. 30；$5\times6=30$（或 $6\times5=30$）\quad 2. 100；3\quad
3. 厘米；米\quad 4. 56\quad 5. 7\quad 6. 3；4\quad 7. 3 个 4；$3\times4=12$\quad
8. $=$\quad 9. 5\quad 10. 4 \\
\textbf{四}：1. ×\quad 2. ×（四四十六只有一个）\quad 3. ×\quad 4. √\quad
5. ×（应是 52）\\
\textbf{五}：1. A\quad 2. A\quad 3. B\quad 4. B\quad 5. A \\
\textbf{六}：1. $8\times5=40$（个）\quad 2. $36\div4=9$（袋）\quad
3. $63\div7=9$\quad 4. $9\times5=45$（元），$45<50$ 够；$50-45=5$（元），还剩 5 元。
\end{ansbox}

\subsection{模拟卷 B（图形、测量与综合）}

\noindent\textbf{一、填一填（每空 1 分，共 16 分）}
\begin{enumerate}[label=\arabic*.,itemsep=3pt]
  \item 填单位：一张床长约 2 \blank[1.4cm]；一块橡皮厚约 1 \blank[1.4cm]。
  \item $1\ \text{米}=$ \blank[1.2cm] 厘米；\ $180\ \text{厘米}=$ \blank[1.2cm] 米
        \blank[1.2cm] 厘米。
  \item $1\ \text{元}=$ \blank[1.2cm] 角；\ $1\ \text{角}=$ \blank[1.2cm] 分；
        $2\ \text{元}5\ \text{角}=$ \blank[1.2cm] 角。
  \item 地图上一般是上 \blank[1cm]、下 \blank[1cm]、左 \blank[1cm]、右 \blank[1cm]。
  \item 小红面向西，她左手边是 \blank[1.2cm]，右手边是 \blank[1.2cm]。
  \item 正方形纸对折后剪一个图形，展开后两边 \blank[2cm]（填“一样”或“不一样”）。
\end{enumerate}

\noindent\textbf{二、量一量、画一画（每题 4 分，共 12 分）}
\begin{enumerate}[label=\arabic*.,itemsep=6pt]
  \item 量一量家里三样东西的长度，写在下面。
        \wls[2]
  \item 画一条长 5 厘米的线段。
        \wls[3]
  \item 画一条长 7 厘米的线段。
        \wls[3]
\end{enumerate}

\noindent\textbf{三、下面哪些是轴对称图形？在括号里打“√”（每题 2 分，共 10 分）}
\begin{center}
\begin{tikzpicture}[scale=0.62]
  \draw[thick] (0,0) circle (0.7);
  \node[below] at (0,-1.0) {圆（\quad）};
  \draw[thick] (2.1,-0.7) rectangle (3.5,0.7);
  \node[below] at (2.8,-1.0) {正方形（\quad）};
  \draw[thick] (4.6,-0.7) -- (6.2,-0.7) -- (5.8,0.7) -- (4.2,0.7) -- cycle;
  \node[below] at (5.2,-1.0) {平行四边形（\quad）};
  \draw[thick] (7.0,-0.7) -- (8.4,-0.7) -- (7.7,0.9) -- cycle;
  \node[below] at (7.7,-1.0) {等腰三角形（\quad）};
  \draw[thick] (9.4,-0.7) -- (10.9,-0.7) -- (10.55,0.7) -- (9.75,0.7) -- cycle;
  \node[below] at (10.15,-1.0) {等腰梯形（\quad）};
\end{tikzpicture}
\end{center}

\noindent\textbf{四、判断（每题 2 分，共 8 分）}
\begin{enumerate}[label=\arabic*.,itemsep=2pt]
  \item 把纸对折后剪出的图形，一定是对称的。\hfill（\quad）
  \item 平移时物体只能横着移动。\hfill（\quad）
  \item 陀螺的轴扎在纸板中心，转起来最稳。\hfill（\quad）
  \item 量长度时，物体一端要对准尺子的 0 刻度。\hfill（\quad）
\end{enumerate}

\noindent\textbf{五、解决问题（每题 6 分，共 24 分）}
\begin{enumerate}[label=\arabic*.,itemsep=8pt]
  \item 一根跳绳长 2 米，2 米等于多少厘米？
        \wls[2]
  \item 妈妈买了 3 支笔，每支 6 元，又买了一个 12 元的文具盒，一共花多少元？
        \wls[3]
  \item 有 50 元，买 6 袋洗衣液，每袋 9 元，钱够吗？
        \wls[3]
  \item 学校在书店东面，邮局在学校东面。从书店到邮局要往哪个方向走？
        \wls[2]
\end{enumerate}

\begin{ansbox}
\textbf{卷 B 答案}\\
\textbf{一}：1. 米；厘米\quad 2. 100；1 米 80 厘米\quad 3. 10；10；25\quad
4. 北、南、西、东\quad 5. 南；北\quad 6. 一样 \\
\textbf{二}：按实际长度作答；第 2、3 题要有端点，并标“5 厘米”“7 厘米”。\\
\textbf{三}：圆 √、正方形 √、平行四边形 ×、等腰三角形 √、等腰梯形 √。\\
\textbf{四}：1. √\quad 2. ×（也能竖着）\quad 3. √\quad 4. √ \\
\textbf{五}：1. $2\ \text{米}=200\ \text{厘米}$\quad
2. $6\times3=18$（元），$18+12=30$（元）\quad
3. $9\times6=54$（元），$54>50$，\textbf{不够}\quad
4. 从书店往东走到学校，再从学校往东走到邮局，\textbf{一直往东走}。
\end{ansbox}

\subsection{考前一天：只做四件事}

\begin{enumerate}
  \item \textbf{看错题本}（10 分钟）：只读错因，不做新题。
  \item \textbf{背口诀}（5 分钟）：跳着抽 20 句，重点看平时总错的。
  \item \textbf{检查文具}：两支削好的铅笔、橡皮、15 cm 直尺
        （\textbf{一定带尺子}，测量题必用）、草稿纸。
  \item \textbf{跟小豆角说一句话}——说具体的，别说“考好点”。比如：
        “上次那个 100 减 48，小豆角自己想到用计数器，那次很棒。”
\end{enumerate}

\subsection{全册最容易错的十个地方}

\begin{center}
\small
\begin{tabularx}{\linewidth}{@{}p{0.7cm}p{5.2cm}X@{}}
\toprule
\textbf{序} & \textbf{易错点} & \textbf{一句话对策} \\
\midrule
1 & 100 减两位数，十位是 0 借不到 & 先上计数器，把“借”做出来 \\
2 & 退位后十位忘了减 1 & 退位点画在十位上，见点先减 \\
3 & 个位不够减时倒着减 & 先借，不许倒着减 \\
4 & 把不同加数的加法改成乘法 & 加数一样才能改 \\
5 & “是几倍”当成“多几个” & 圈一圈，圈出几组就是几倍 \\
6 & 求“是几倍”用了乘法 & 看里面有几个 $\to$ 用除法 \\
7 & 括号里最大填几填大了 & 填完必须乘一遍 \\
8 & 量长度从尺子边缘开始 & 手指点住 0 \\
9 & 坏尺子不会用 & 起点 + 长度 = 终点 \\
10 & 应用题单位、答语漏写 & 写完回头读一遍问题 \\
\bottomrule
\end{tabularx}
\end{center}

\noindent\textbf{这十条要是都能对号、能说清，二年级上册就没有真正的难点了。}
